% Level 3, Week 8 -- Data Detectives
\lesson{8}{Data Detectives}{Data tells stories when you clean it, count it and chart it.}
\lessoninfo{Concepts: CSV files, csv.DictReader, cleaning, counting with dicts, averages, bar charts}{Tools: Thonny, turtle, csv}{CSTA: 2-DA-08, 3A-DA-11, 3A-DA-12, 3A-IC-29}

\begin{objectives}
  \item The student can read a CSV file with \texttt{csv.DictReader} and clean it: strip spaces, skip empty rows, turn text into numbers. (L3.G6)
  \item The student can count values with a dictionary and work out an average that skips missing data. (L3.G3, L3.G6)
  \item The student can draw a bar chart with turtle where every bar is labelled with text, not colour alone. (L3.G6)
  \item The student can explain why surveys need consent and how to anonymise answers before sharing. (L3.G9)
\end{objectives}

\begin{materials}
  \item Print: workbook chapter 8 and homework 8 (one per student).
  \item On each Mac: in Thonny open \texttt{3-advanced/excercises/week08/ex-survey-cleaner/starter.py} and \texttt{favourite\_games.csv} next to it (in \menu{File} $\rightarrow$ \menu{Open} choose \enquote{All files}).
  \item Unplugged: one sticky note per person (including you), a large sheet of paper, a marker.
  \item Prep (10 min): read this plan; run \texttt{ex-bar-chart/solution.py} once to see the finished chart.
\end{materials}

\begin{timeplan}
  \timeslot{0--8}{Warm-up}{Unplugged paper survey with consent: anonymous sticky notes $\rightarrow$ tally $\rightarrow$ hand-drawn bar chart.}
  \timeslot{8--18}{Concept}{What a CSV file is; \texttt{DictReader} gives one dictionary per row; why data must be cleaned; consent and anonymising.}
  \timeslot{18--28}{Guided demo}{Type-along: Exercise 8.1 survey cleaner (19 messy rows $\rightarrow$ 17 clean ones).}
  \timeslot{28--52}{Hands-on build}{Pairs: Exercise 8.2 stats (about 10 min), then 8.3 bar chart (about 10 min); swap every 8 min. Bug hunt. Last 4 min: free-build chart.}
  \timeslot{52--60}{Wrap-up}{What story does our data tell? Exit ticket, hand out homework 8.}
\end{timeplan}

\section{Warm-up (0--8 min)}
\begin{unplugged}[The sticky-note survey]
Say: \enquote{I would like to run a survey: \emph{What is your favourite game?} Taking part is your choice -- you can pass, and you do not write your name.} (This is \textbf{consent} and \textbf{anonymity}; you will name them in a moment.) Everyone, you included, writes one game on a sticky note, folds it and drops it in a pile. Shuffle. A volunteer reads them out; another makes a tally on the big sheet. Notice together: \enquote{Did someone write \emph{minecraft} and someone else \emph{Minecraft!}? Are they the same answer?} Agree on one spelling -- you just \textbf{cleaned} data. Draw a quick bar chart with a label under every bar.
\end{unplugged}

\section{Concept (8--18 min)}
\begin{teacherscript}
\enquote{Real data detectives get answers from hundreds of people, so they keep them in a file. A \textbf{CSV} file -- comma-separated values -- is a table written as plain text: one line per row, commas between columns. The first line is the \textbf{header}: the column names.

Python's \py|csv.DictReader| reads each row into a \textbf{dictionary}: column name $\rightarrow$ value. But every value arrives as \emph{text}: \py|"5"|, not \py|5|. And real data is messy: extra spaces, MiXeD capitals, empty rows. So the first job is always \textbf{cleaning}. Then we \textbf{count}, \textbf{average} and \textbf{chart}.}
\end{teacherscript}

Show \texttt{favourite\_games.csv} in Thonny and hunt for mess together: \texttt{ block craft } (spaces), \texttt{BLOCK CRAFT} (capitals), row 5 and an all-commas row (empty), row 7 (missing hours).

\begin{analogy}[Washing vegetables]
Data is like vegetables from the garden: you would not cook them with the mud still on. Cleaning comes first -- strip off the spaces, throw away the empty rows, cut text into numbers -- then you cook (count, average, chart).
\end{analogy}

\begin{teachernote}[The CS idea, for you]
\enquote{Garbage in, garbage out}: counting messy data gives confidently wrong answers. Missing values become \texttt{None} and are skipped in averages -- counting them as 0 would drag the average down unfairly.
\end{teachernote}

\section{Guided demo (18--28 min)}
\begin{workedexample}[Exercise 8.1: survey cleaner (type-along)]
\begin{enumerate}[itemsep=1pt]
  \item Run \texttt{starter.py}. It says 19 answers and prints \texttt{block craft} with spaces. \enquote{Our detective is being fooled!}
  \item TODO 1, \texttt{to\_number}: strip; empty text gives \py|None|; otherwise \py|int(text)|.
  \item TODO 2, \texttt{is\_empty}: \py|return row["favourite_game"].strip() == ""|.
  \item TODO 3: add \py|.title()| so all spellings become \texttt{Block Craft}. Run: \texttt{Clean answers: 17}.
\end{enumerate}
\end{workedexample}
\pyfile[firstline=25,lastline=45]{../excercises/week08/ex-survey-cleaner/solution.py}

\section{Hands-on build (28--52 min)}
\begin{pairroles}[Detective and recorder]
The \textbf{driver} types. The \textbf{navigator} is the \enquote{recorder}: before each Run they write on paper what answer they \emph{expect} (e.g. Block Craft has 5 votes), then check. Swap at 8 minutes. With 1 student: you are the recorder. With 3: add a \textbf{checker} who runs \texttt{check\_my\_work.py}; rotate. With 4: two pairs; at the end each pair explains one function to the other.
\end{pairroles}

\textbf{Exercise 8.2} (about 10 min): \texttt{count\_by}, \texttt{average}, \texttt{most\_popular}. \textbf{Exercise 8.3} (about 10 min): scale the bars and write a name and a number on every bar.

\begin{bughunt}[Numbers that are really words]
In the Thonny shell type \py|hours = ["5", "3", "2"]| then \py|max(hours), sum(hours)|. \texttt{max} says \py|"5"| (looks fine!), but \texttt{sum} crashes: \texttt{TypeError: unsupported operand type(s) for +: 'int' and 'str'}. Try \py|max(["10", "9"])| -- it says \texttt{"9"}, because text is compared letter by letter. Fix: convert with \py|int()| while cleaning. Then the second bug: \py|int("")| crashes on an empty answer -- that is why \texttt{to\_number} returns \py|None| and \texttt{numbers\_in} skips it.
\end{bughunt}

\begin{freebuild}[Your chart, your question (last 4 minutes)]
Chart a different column: change \texttt{count\_games} to count \texttt{genre} or \texttt{age\_group}. Or add a y-axis with a tick for every vote, or sort the bars alphabetically.
\end{freebuild}

\section{Wrap-up (52--60 min)}
Look at the chart together. Ask: \enquote{What story does it tell? What can it \emph{not} tell us?} (Only 17 people answered; \enquote{adventure} has 5.0 stars from just 2 answers -- small samples can mislead.)

\begin{formativecheck}[Exit ticket]
\textbf{Look for:} students say \enquote{clean before counting}, and every bar has a label.
\textbf{Question:} \enquote{Answers for hours: 4, 2, (blank), 6. What is the average, and what did you do with the blank?}
\textbf{Answer:} 4 -- skip the blank: (4 + 2 + 6) / 3. Counting it as 0 would give 3, which is unfair.
\end{formativecheck}

\begin{reflection}
What was the messiest thing in the data? Why do we write the name on each bar instead of only using colours? Would you answer a survey that asked for your full name? Why or why not?
\end{reflection}

\begin{safety}[Consent and anonymising survey data]
Say: \enquote{Our survey asked first, let people pass, and had no names. That is \textbf{consent} and \textbf{anonymity}. Even without names, some answers together can point to one person -- \emph{the only 13-year-old in the chess club}. Before sharing data: keep only the columns you need, replace names with IDs like R1, R2, and never collect addresses, schools, birthdays or photos.} Homework 8 practises this.
\end{safety}

Hand out homework 8 (survey privacy: anonymise a CSV file).

\section{Notes for the teacher}
\begin{commonmistakes}
  \item \textbf{Forgetting \texttt{import csv}} or opening the file from the wrong folder: \texttt{FileNotFoundError}. Open \texttt{starter.py} from the exercise folder so the CSV sits next to it.
  \item \textbf{Column name typos}: \py|row["game"]| instead of \py|row["favourite_game"]| gives \texttt{KeyError}. Copy names from the header line.
  \item \textbf{Counting without \texttt{.get}}: \py|counts[value] += 1| crashes the first time (\texttt{KeyError}). Use \py|counts.get(value, 0) + 1|.
  \item \textbf{Bars not scaled}: dividing by the wrong number, or integer maths. \py|value / biggest * MAX_HEIGHT| works.
\end{commonmistakes}

\begin{differentiation}
\supportpath Do 8.1 and 8.2 together as a group; skip 8.3 or run the solution and change only the colours and title. Trace \texttt{count\_by} on paper with a 4-row table.
\extensionpath Work out the average stars \emph{per genre} (the solution's main block shows how), or write \texttt{save\_report()} that writes the results to \texttt{report.txt}.
\end{differentiation}

\begin{stretchchallenge}[Two-column detective]
Make a grouped chart: for each game, two bars side by side -- one for the 10--11 group, one for the 12--13 group -- each with a text label. Which age group likes puzzle games more?
\end{stretchchallenge}

\begin{teachernote}[If you are short on time]
Shorten the warm-up to the tally only (no paper chart). Do 8.1 as the demo and 8.3 in pairs (it has its own loader); 8.2 can move to Buffer B2.
\end{teachernote}
